% main.tex
\documentclass[11pt,letterpaper]{article}

% --- 基本套件 ---
\usepackage[utf8]{inputenc}
\usepackage[T1]{fontenc}
\usepackage{lmodern}
\usepackage[margin=1in]{geometry}
\usepackage{amsmath,amssymb,amsthm}
\usepackage{mathtools}
\usepackage{graphicx}
\usepackage{booktabs}
\usepackage{multirow}
\usepackage{siunitx}
\usepackage{xcolor}
\usepackage{hyperref}
\usepackage{cleveref}
\usepackage{authblk}
\usepackage{caption}
\usepackage{subcaption}
\usepackage{algorithm}
\usepackage{algpseudocode}
\usepackage{enumitem}
\usepackage{threeparttable}
\usepackage{float}

% --- 超連結設定 ---
\hypersetup{
    colorlinks=true,
    linkcolor=blue!60!black,
    citecolor=green!50!black,
    urlcolor=blue!70!black,
}

% --- 定理環境 ---
\newtheorem{theorem}{Theorem}
\newtheorem{proposition}{Proposition}
\newtheorem{lemma}{Lemma}
\newtheorem{corollary}{Corollary}
\theoremstyle{definition}
\newtheorem{definition}{Definition}
\newtheorem{remark}{Remark}

% --- 自定義命令 ---
\newcommand{\R}{\mathbb{R}}
\newcommand{\N}{\mathbb{N}}
\newcommand{\ZSP}{\textsc{ZSP}\xspace}
\newcommand{\MOT}{\textsc{MOT}\xspace}
\newcommand{\abnqss}{\textsc{ABN-QSS}\xspace}
\newcommand{\facs}{\textsc{FACS}\xspace}
\newcommand{\proj}{\mathcal{P}}
\newcommand{\nullsp}{\mathrm{Null}}
\newcommand{\rank}{\mathrm{rank}}
\newcommand{\diag}{\mathrm{diag}}
\newcommand{\Var}{\mathrm{Var}}
\newcommand{\Mmat}{\mathbf{M}}
\newcommand{\Cvec}{\mathbf{C}}
\newcommand{\xvec}{\mathbf{x}}
\newcommand{\avec}{\mathbf{a}}
\newcommand{\zvec}{\mathbf{z}}
\newcommand{\Vmat}{\mathbf{V}}
\newcommand{\evec}{\mathbf{e}}

\usepackage{xspace}

% --- 標題 ---
\title{
    \bfseries The Boundaries of Physics-Native Computing:\\
    A Systematic Negative Result on Zero-Space Projection\\
    for Multi-Object Tracking and Cooperative Perception
}

\author[1]{Lien-Hsing Chang}
\author[2]{ABN-QSS Research Team}
\affil[1]{Independent Researcher}
\affil[2]{\texttt{abn-qss@research.org}}

\date{September 2026}

\begin{document}

\maketitle

\begin{abstract}
Zero-Space Projection (\ZSP) is a core mechanism of the 
\abnqss v3.0 physics-native computing architecture, whose digital 
twin has demonstrated projection linearity $R^2 = 1.0$ and a 
null-space dimension formula $K = N(N-1)/2$. This paper 
systematically evaluates \ZSP on three perception tasks: 
single-frame multi-object tracking (\MOT), multi-frame \MOT, and 
cooperative perception. Our findings are threefold: 
(i) on \textbf{single-frame \MOT}, \ZSP reaches $99.96\%$ of the 
accuracy of the Hungarian algorithm but is $20{,}000\times$ slower; 
(ii) on \textbf{multi-frame \MOT}, a paired $t$-test across 30 seeds 
shows \ZSP is $6.7\%$ worse than Hungarian on average 
($p = 0.0001$, win rate $0\%$); 
(iii) on \textbf{cooperative perception}, all three methods 
(centralized Hungarian, distributed Hungarian, and distributed \ZSP) 
achieve $100\%$ accuracy, but distributed Hungarian is fastest with 
zero communication overhead. 

We analyze the root cause: linear assignment problems possess a 
\emph{dual structure} that enables the Hungarian algorithm to achieve 
exact optimality in polynomial time, whereas \ZSP is an iterative 
approximation that cannot, by construction, outperform the exact 
algorithm. We conclude that the value of \ZSP---and of physics-native 
computing more broadly---lies not in \emph{replacing} problems already 
solved by exact algorithms, but in \emph{solving problems without good 
algorithms}: material simulation, physics-layer cryptography, 
high-precision sensing, and variational field solving. This work 
provides an explicit boundary for physics-native computing paradigms.
\end{abstract}

\tableofcontents
\newpage

\section{Introduction}
\label{sec:intro}

\subsection{Motivation: Physics-Native Computing}

Since 1945, the von Neumann architecture has defined digital 
computation: programs are stored in memory and executed by a 
sequential fetch-decode-execute cycle. As Moore's law slows and the 
memory wall tightens, this separation imposes two fundamental limits: 
(a) \emph{energy waste}: moving data costs $100\times$ more than 
computing on it~\cite{horowitz2014computing}; and 
(b) \emph{the clock bottleneck}: complex iterative problems 
(such as Einstein field equations, fluid dynamics, and quantum 
many-body systems) suffer severe latency from sequential iteration 
and ECC verification overhead.

Physics-native computing proposes a different paradigm: rather than 
\emph{simulating} mathematics in hardware, let the hardware 
\emph{be} the mathematics. The intrinsic physical laws of analog 
circuits---Ohm's law, Kirchhoff's laws, capacitor charge dynamics---can 
directly realize variational relaxation and solve constraint-satisfaction 
problems without clocked iteration~\cite{adamatzky2019physics}.

\subsection{Zero-Space Projection: A Brief Recap}

\abnqss v3.0 is a representative physics-native architecture based on 
a silicon-photonics MZI array (L0 layer) coupled with an analog OTA-C 
array (L1 ``Shenqu engine''). Its core theoretical claim is the 
\emph{zero-space projection lock theorem}:
\begin{equation}
    \Cvec_{\mathrm{final}} = \proj \cdot (\Cvec_{\mathrm{base}} 
    - \delta \cdot \alpha \cdot \evec_0),
    \label{eq:projection}
\end{equation}
where $\proj = \Vmat \Vmat^\top$ is the projection matrix onto the 
null space of the constraint matrix $\Mmat$, $\Vmat$ is a null-space 
basis, $\delta$ is a physical perturbation, and $\alpha$ is the 
physical coupling constant.

The digital twin of \abnqss v3.0 has been validated for 
$N = 3, \dots, 6$, showing:
\begin{itemize}[leftmargin=*]
    \item Projection linearity $R^2 = 1.000000$ (exact);
    \item Null-space dimension $K = N(N-1)/2$ (exact);
    \item Physical coupling slope $\alpha \in [0.0075, 0.0115]$, 
          increasing with $N$.
\end{itemize}

\subsection{Motivating the Present Study}

Given the strong theoretical foundation of \ZSP, a natural question 
arises: \emph{where should \ZSP be applied?} 

Multi-object tracking (\MOT) and cooperative perception are natural 
candidate applications, because their core computational step---data 
association---can be cast as a constraint-satisfaction problem. Several 
prior works have advocated for physics-based or analog approaches to 
perception~\cite{...}. 

This paper reports the outcome of a six-month systematic study 
evaluating \ZSP on three perception tasks. \textbf{Our results are 
negative}: in all three tasks, the classical Hungarian algorithm 
is either strictly better (in accuracy, speed, or communication 
cost) or equal. 

\subsection{Contributions}

Our contributions are:
\begin{enumerate}[leftmargin=*]
    \item \textbf{A complete experimental framework} spanning 
          digital-twin verification to real benchmarks, with 
          all code and data released open-source.
    \item \textbf{Systematic negative results} on three perception 
          tasks: single-frame \MOT, multi-frame \MOT, and 
          cooperative perception.
    \item \textbf{A root-cause analysis}: linear assignment problems 
          possess a dual structure that enables exact polynomial 
          algorithms, which \ZSP's iterative nature cannot beat.
    \item \textbf{An explicit applicability boundary}: \ZSP should 
          be applied to problems without good algorithms, not to 
          problems already solved by exact methods.
\end{enumerate}

\subsection{Paper Organization}

\Cref{sec:related} reviews related work. 
\Cref{sec:method} formalizes \ZSP for the three perception tasks. 
\Cref{sec:experiments} presents the experimental results. 
\Cref{sec:analysis} analyzes why \ZSP fails. 
\Cref{sec:discussion} discusses the applicability boundary. 
\Cref{sec:conclusion} concludes.
\section{Related Work}
\label{sec:related}

\subsection{Data Association in Multi-Object Tracking}

Data association in \MOT is a classical linear assignment problem. 
Traditional methods include:
\begin{itemize}[leftmargin=*]
    \item \textbf{Hungarian algorithm}~\cite{kuhn1955hungarian}: 
          an $O(n^3)$ exact polynomial algorithm.
    \item \textbf{MIP methods}~\cite{jiang2015multiple}: integer 
          programming, exact at small scales but slow at large scales.
    \item \textbf{Greedy matching}: $O(n \log n)$ but non-optimal.
    \item \textbf{Deep learning approaches}~\cite{zhang2022bytetrack,
          cao2023observation}: learn similarity functions, but 
          association itself still uses Hungarian.
\end{itemize}

A key observation is that the \emph{combinatorial} structure of \MOT 
data association is \emph{linear}. Its dual form admits exact polynomial 
solution. This is the theoretical reason for Hungarian's dominance.

\subsection{Cooperative Perception}

Cooperative perception~\cite{chen2019cooper} enables multiple vehicles 
or edge devices to share sensor data and overcome single-vehicle field-
of-view limitations. Representative methods include:
\begin{itemize}[leftmargin=*]
    \item \textbf{V2X-ViT}~\cite{xu2022v2x}: Transformer-based 
          cooperative perception.
    \item \textbf{CoBEVT}~\cite{xu2022cobevt}: bird's-eye view 
          cooperative perception.
    \item \textbf{AttFuse}~\cite{xu2022opv2v}: attention-based 
          intermediate fusion.
\end{itemize}
Again, the association step in these pipelines remains Hungarian.

\subsection{Physics-Native Computing and \abnqss}

\abnqss v3.0~\cite{chang2026abnqss} is a representative 
physics-native computing architecture with:
\begin{itemize}[leftmargin=*]
    \item \textbf{L0 photonic layer}: silicon-photonics MZI 
          interferometer array for matrix multiplication and 
          null-space projection.
    \item \textbf{L1 analog layer}: OTA-C ``Shenqu engine'' for 
          nonlinear feedback control.
    \item \textbf{L2 digital layer}: FPGA/CMOS for timing and 
          buffering.
    \item \textbf{L5 cognitive layer}: external host/SoC for 
          macro task scheduling.
\end{itemize}
The digital twin of \abnqss v3.0 has been validated for $N=3$ to 
$N=6$ via Python simulation~\cite{chang2026n4}.

\subsection{Hungarian Algorithm and the Dual Structure}

Linear assignment is:
\begin{equation}
    \min_{x \in \{0,1\}^{T \times D}} \sum_{i,j} c_{ij} x_{ij}
    \quad \text{s.t.} \quad 
    \sum_i x_{ij} = 1,\; \sum_j x_{ij} = 1.
\end{equation}
Its Lagrangian dual is:
\begin{equation}
    \max_{u, v} \sum_i u_i + \sum_j v_j
    \quad \text{s.t.} \quad u_i + v_j \le c_{ij}.
\end{equation}
The Hungarian algorithm exploits this dual structure and terminates 
in $O(n^3)$ time with an exact certificate of optimality.

\textbf{Theoretical consequence}: any iterative approximation 
algorithm (including \ZSP) cannot improve on Hungarian for this 
problem class, because the dual certificate is already exact.
\section{Methodology}
\label{sec:method}

\subsection{Null-Space Projection Framework}

Given a constraint matrix $\Mmat \in \R^{m \times n}$, the null 
space is:
\begin{equation}
    \nullsp(\Mmat) = \{\Cvec \in \R^n : \Mmat \Cvec = 0\}.
\end{equation}
A null-space basis $\Vmat \in \R^{n \times K}$ satisfies $\Mmat \Vmat = 0$. 
Any feasible solution can be parameterized as:
\begin{equation}
    \Cvec = \Vmat \avec + \Cvec_0, \quad \avec \in \R^K,
\end{equation}
where $\Cvec_0$ is a particular solution and $K = n - \rank(\Mmat)$.

The projection matrix is $\proj = \Vmat \Vmat^\top$.

\subsection{Formalization for MOT Association}

Consider $T$ tracks and $D$ detections. Let $x_{ij} \in [0,1]$ denote 
the association probability. The constraints are:
\begin{align}
    \sum_{i=1}^{T} x_{ij} + s_j &= 1, \quad j = 1, \dots, D, 
        \label{eq:det-uniqueness}\\
    \sum_{j=1}^{D} x_{ij} + t_i &= 1, \quad i = 1, \dots, T,
        \label{eq:traj-uniqueness}
\end{align}
where $s_j, t_i \ge 0$ are slack variables.

\paragraph{Homogenization.} Define 
$\Cvec = [\xvec^\top, \mathbf{s}^\top, \mathbf{t}^\top, 1]^\top$ 
and $\Mmat = [\mathbf{A}, -\mathbf{b}]$ where 
$\mathbf{A} \Cvec_{1:n} = \mathbf{b}$ encodes 
\cref{eq:det-uniqueness,eq:traj-uniqueness}. Then the constraint 
becomes $\Mmat \Cvec = 0$.

\subsection{Multi-Frame Extension with Soft Continuity}

For $K$ frames, we stack the association matrices 
$\xvec^{(1)}, \dots, \xvec^{(K)}$ into a single variable. In addition 
to per-frame constraints, we introduce a \emph{soft continuity} 
penalty:
\begin{equation}
    J_{\text{cont}} = \beta \sum_{t=1}^{K-1} 
    \left\| \mathbf{r}^{(t+1)} - \mathbf{r}^{(t)} \right\|^2,
\end{equation}
where $r_i^{(t)} = \sum_j x_{ij}^{(t)}$ is the row-sum for track $i$ 
at frame $t$, and $\beta$ controls the continuity weight.

The full objective is:
\begin{equation}
    J = \underbrace{-\sum_t \sum_{i,j} p_{ij}^{(t)} x_{ij}^{(t)}}_{
        \text{similarity reward}}
    + \lambda \|\xvec\|^2 
    + \mu (\|\mathbf{s}\|^2 + \|\mathbf{t}\|^2)
    + \beta \|J_{\text{cont}}\|.
\end{equation}

\subsection{Cooperative Perception Extension}

For $M$ vehicles, each with $T$ local tracks and $D$ local detections, 
we solve a distributed null-space problem. Each vehicle solves:
\begin{equation}
    \min_{\avec_m} J_m(\Vmat_m \avec_m),
\end{equation}
and exchanges only the association vectors of boundary detections:
\begin{equation}
    \mathbf{b}_m = \{ x_{ij}^{(m)} : (i,j) \in \mathcal{B}_m \},
\end{equation}
where $\mathcal{B}_m$ is the boundary set (detections shared with 
neighbors).

\paragraph{Communication cost.} Each vehicle broadcasts 
$T \times 8$ bytes per boundary round, totaling 
$M \times T \times 8 \times R$ bytes for $R$ rounds.

\subsection{Projected Gradient Solver}

We solve the problem via projected gradient descent on $\avec$:
\begin{equation}
    \avec_{k+1} = \avec_k - \eta \Vmat^\top \nabla_\Cvec J(\Vmat \avec_k) 
    - \gamma (D - D_{\text{th}}) \frac{\avec_k - \bar{\avec}_k}{\|\avec_k - \bar{\avec}_k\|},
\end{equation}
where:
\begin{itemize}[leftmargin=*]
    \item $\eta$ is the learning rate;
    \item $\gamma$ is the symmetry recovery gain;
    \item $D = \Var(\xvec) / \sigma_{\max}^2$ is the mismatch;
    \item $D_{\text{th}}$ is the mismatch threshold.
\end{itemize}

After convergence, we binarize via the Hungarian algorithm on the 
continuous solution.

\subsection{Convergence Guarantee}

The loss $J$ is convex and $L$-smooth on the null-space parameterization. 
By standard projected gradient theory, the update converges 
monotonically to the global optimum for $\eta < 2/L$. \Cref{app:convergence} 
provides details.
\section{Experiments}
\label{sec:experiments}

\subsection{Single-Frame MOT Benchmark}

\paragraph{Setup.} Random similarity matrices 
$p \in [0,1]^{T \times D}$ with $T = D \in \{10, 20, 30, 50\}$.

\begin{table}[t]
    \centering
    \caption{Single-frame MOT: solution quality and computation time.}
    \label{tab:single-frame}
    \begin{threeparttable}
    \begin{tabular}{l S[table-format=-2.2] S[table-format=-2.2] 
                    S[table-format=-2.2] }
        \toprule
        \multirow{2}{*}{Size $T=D$} & 
        \multicolumn{3}{c}{Assignment cost} \\
        \cmidrule(lr){2-4}
        & {Hungarian} & {Greedy} & {\ZSP} \\
        \midrule
        10 & -8.60 & -7.95 & -8.59 \\
        20 & -18.46 & -17.52 & -18.46 \\
        30 & -28.51 & -27.07 & -28.50 \\
        50 & -48.39 & -46.77 & -48.37 \\
        \bottomrule
    \end{tabular}
    \end{threeparttable}
\end{table}

\begin{table}[t]
    \centering
    \caption{Single-frame MOT: computation time (ms).}
    \label{tab:single-frame-time}
    \begin{tabular}{l S[table-format=0.2] S[table-format=4.2] 
                    S[table-format=5.0] }
        \toprule
        Size $T=D$ & {Hungarian} & {\ZSP} & {Ratio} \\
        \midrule
        10 & 0.01 & 34.56 & 3456 \\
        20 & 0.01 & 65.14 & 6514 \\
        30 & 0.03 & 222.27 & 7409 \\
        50 & 0.07 & 1498.87 & 21412 \\
        \bottomrule
    \end{tabular}
\end{table}

\paragraph{Result.} \ZSP achieves $99.96\%$ of Hungarian's solution 
quality but is $3456$--$21412\times$ slower 
(\cref{tab:single-frame,tab:single-frame-time}). 
See \cref{fig:single-frame} for visual comparison.

\begin{figure}[t]
    \centering
    \includegraphics[width=\linewidth]{figures/fig2_single_frame.pdf}
    \caption{Single-frame MOT: (a) computation time on log scale; 
    (b) solution quality. Hungarian is both faster and (marginally) 
    more accurate than \ZSP.}
    \label{fig:single-frame}
\end{figure}

\subsection{Multi-Frame MOT Benchmark}

\paragraph{Setup.} Synthetic data with $K$ frames, $T$ tracks, $D$ 
detections, incorporating miss detections ($25\%$), false positives 
($15\%$), and occlusion ($20\%$). We run 30 seeds per configuration.

\begin{table}[t]
    \centering
    \caption{Multi-frame MOT: statistical comparison across 30 seeds.}
    \label{tab:multiframe}
    \begin{tabular}{l S[table-format=1.3] S[table-format=1.3] 
                    S[table-format=+1.3] }
        \toprule
        Metric & {Hungarian} & {\ZSP} & {Difference} \\
        \midrule
        Mean accuracy & 0.909 & 0.842 & -0.067 \\
        Std. deviation & 0.122 & 0.129 & --- \\
        $t$-statistic & \multicolumn{2}{c}{$-4.580$} & \\
        $p$-value & \multicolumn{2}{c}{$0.0001$} & \\
        Win rate & \multicolumn{2}{c}{$0\%$} & \\
        \bottomrule
    \end{tabular}
\end{table}

\paragraph{Result.} \ZSP is \emph{statistically significantly worse} 
than Hungarian ($-6.7\%$, $p=0.0001$), with zero wins across 
30 seeds (\cref{tab:multiframe,fig:multiframe}).

\begin{figure}[t]
    \centering
    \includegraphics[width=\linewidth]{figures/fig3_multiframe.pdf}
    \caption{Multi-frame MOT: (a) paired accuracy scatter; 
    (b) difference distribution; (c) CDF. \ZSP is systematically 
    worse than Hungarian.}
    \label{fig:multiframe}
\end{figure}

\subsection{Cooperative Perception Benchmark}

\paragraph{Setup.} $M$ vehicles, each with $T$ tracks and $D$ 
detections, sharing boundary data.

\begin{table}[t]
    \centering
    \caption{Cooperative perception: accuracy, time, communication.}
    \label{tab:cooperative}
    \begin{tabular}{l l S[table-format=1.3] S[table-format=3.2] 
                    S[table-format=4.0]}
        \toprule
        Config & Method & {Accuracy} & {Time (ms)} & {Comm (B)} \\
        \midrule
        \multirow{3}{*}{$M{=}2,T{=}6,D{=}8$} 
            & Central.\ Hung. & 1.000 & 21.71 & --- \\
            & Distrib.\ Hung. & 1.000 & 0.03 & 0 \\
            & Distrib.\ \ZSP & 1.000 & 15.99 & 480 \\
        \midrule
        \multirow{3}{*}{$M{=}3,T{=}8,D{=}10$} 
            & Central.\ Hung. & 1.000 & 0.82 & --- \\
            & Distrib.\ Hung. & 1.000 & 0.04 & 0 \\
            & Distrib.\ \ZSP & 1.000 & 26.22 & 960 \\
        \midrule
        \multirow{3}{*}{$M{=}4,T{=}10,D{=}12$} 
            & Central.\ Hung. & 1.000 & 1.56 & --- \\
            & Distrib.\ Hung. & 1.000 & 0.05 & 0 \\
            & Distrib.\ \ZSP & 1.000 & 71.23 & 1600 \\
        \bottomrule
    \end{tabular}
\end{table}

\paragraph{Result.} All methods achieve $100\%$ accuracy, but 
distributed Hungarian is fastest and requires zero communication 
(\cref{tab:cooperative,fig:cooperative}).

\begin{figure}[t]
    \centering
    \includegraphics[width=\linewidth]{figures/fig4_cooperative.pdf}
    \caption{Cooperative perception: (a) computation time on log 
    scale; (b) communication overhead. Distributed Hungarian 
    dominates on both metrics.}
    \label{fig:cooperative}
\end{figure}

\subsection{Noise Resilience}

\begin{figure}[t]
    \centering
    \includegraphics[width=0.75\linewidth]{figures/fig5_noise.pdf}
    \caption{Noise resilience of \ZSP solver: synchronization 
    fitness vs.\ noise intensity $\sigma$. Safe operation up to 
    $\sigma = 0.03$; graceful degradation zone from $0.03$ to $0.05$.}
    \label{fig:noise}
\end{figure}

\subsection{Digital-Twin Verification}

\begin{figure}[t]
    \centering
    \includegraphics[width=\linewidth]{figures/fig1_digital_twin.pdf}
    \caption{Digital-twin verification of \abnqss v3.0: 
    (a) projection linearity $R^2 = 1.000000$; 
    (b) null-space dimension $K = N(N-1)/2$. Both match theory exactly.}
    \label{fig:digital-twin}
\end{figure}
\section{Analysis: Why Does \ZSP Fail?}
\label{sec:analysis}

\subsection{The Dual Structure of Linear Assignment}

The linear assignment problem admits a Lagrangian dual:
\begin{equation}
    \min_{x} c^\top x \quad \text{s.t.} \quad \mathbf{1}^\top x = 1,\; 
    x \mathbf{1} = 1,\; x \ge 0
\end{equation}
has dual:
\begin{equation}
    \max_{u,v} \mathbf{1}^\top u + \mathbf{1}^\top v 
    \quad \text{s.t.} \quad u_i + v_j \le c_{ij}.
\end{equation}

By LP duality, both have the same optimal value. The Hungarian 
algorithm constructs an optimal dual certificate in $O(n^3)$ time. 
This certificate \emph{proves} optimality: no iterative method can 
produce a better solution.

\subsection{Iterative Approximation's Intrinsic Disadvantage}

\begin{table}[t]
    \centering
    \caption{Hungarian vs.\ \ZSP: fundamental properties.}
    \label{tab:comparison}
    \begin{tabular}{l c c}
        \toprule
        Property & Hungarian & \ZSP \\
        \midrule
        Convergence & One-shot & Iterative (hundreds of steps) \\
        Optimality & Exact & Approximate \\
        Complexity & $O(n^3)$ & $O(n^2 K \cdot \text{iter})$ \\
        Dual exploitation & \checkmark & $\times$ \\
        Optimality certificate & \checkmark & $\times$ \\
        \bottomrule
    \end{tabular}
\end{table}

As \cref{tab:comparison} shows, the theoretical gap is 
\emph{structural}, not merely empirical.

\subsection{The Ill-Posed Relaxation}

\MOT data association is a \emph{combinatorial} problem: 
$x_{ij} \in \{0,1\}$. \ZSP relaxes it to $x_{ij} \in [0,1]$ and 
rounds the solution. The relaxation gap is bounded by the integrality 
gap, which is zero for assignment polytopes (the LP relaxation is 
tight). This means:

\begin{proposition}
    The LP relaxation of the linear assignment problem has the same 
    optimal value as the integer problem.
\end{proposition}

\begin{corollary}
    The rounding step in \ZSP cannot recover any quality lost during 
    relaxation, because the relaxation is already tight. The remaining 
    gap is entirely due to \emph{iterative non-convergence}, not to 
    relaxation error.
\end{corollary}

This explains why \ZSP's solution quality ($99.96\%$) is close to 
Hungarian's but never exceeds it.

\subsection{Scaling Analysis}

\begin{figure}[t]
    \centering
    \includegraphics[width=0.7\linewidth]{figures/fig7_scaling.pdf}
    \caption{Scaling behavior: Hungarian is $O(n^{1.26})$ in our 
    measurements (below theoretical $O(n^3)$ due to small sizes); 
    \ZSP is $O(n^{2.33})$.}
    \label{fig:scaling}
\end{figure}

Empirical fits give:
\begin{align}
    \text{Hungarian} &: O(n^{1.26}), \\
    \text{\ZSP} &: O(n^{2.33}).
\end{align}
At $T=D=50$, the gap is already $20{,}000\times$; 
at $T=D=100$, extrapolation predicts $10^6\times$. 
The scaling gap is not a constant but grows without bound.

\subsection{Convergence Behavior}

\begin{figure}[t]
    \centering
    \includegraphics[width=\linewidth]{figures/fig6_convergence.pdf}
    \caption{\ZSP convergence: (a) loss trajectory with initial 
    oscillation; (b) gradient norm decay; (c) mismatch evolution.}
    \label{fig:convergence}
\end{figure}

\Cref{fig:convergence} shows \ZSP's typical convergence: 
initial oscillation due to non-convex geometry, mid-phase rapid 
descent, late-phase locking at the optimum. The 500-iteration 
budget is often exhausted.
\section{Discussion}
\label{sec:discussion}

\subsection{The Applicability Boundary of \ZSP}

\begin{figure}[t]
    \centering
    \includegraphics[width=0.85\linewidth]{figures/fig8_boundary.pdf}
    \caption{Applicability boundary of \ZSP. Problems with dual 
    structure (left) are best solved by exact polynomial algorithms. 
    Problems without such structure (right) are the natural domain 
    of \ZSP.}
    \label{fig:boundary}
\end{figure}

\begin{table}[t]
    \centering
    \caption{Where \ZSP helps and where it doesn't.}
    \label{tab:applicability}
    \begin{tabular}{l c c}
        \toprule
        Problem class & Dual structure? & \ZSP suitable? \\
        \midrule
        Linear assignment (\MOT, cooperative) & Yes & No \\
        Non-convex constraint satisfaction & No & Yes \\
        Physics field simulation & No & Yes \\
        High-dim non-linear optimization & No & Yes \\
        \bottomrule
    \end{tabular}
\end{table}

As \cref{tab:applicability} shows, \ZSP is well-suited for problems 
that lack a dual structure and thus lack exact polynomial algorithms.

\subsection{Implications for Physics-Native Computing}

The original application targets of \abnqss v3.0 are:
\begin{itemize}[leftmargin=*]
    \item \textbf{Material simulation}: entanglement fingerprint 
          analysis via Hadamard grid projection.
    \item \textbf{Secure communication}: physical-layer error 
          correction via $M C = 0$ constraints.
    \item \textbf{High-precision sensing}: molecular detection via 
          locked linearity.
    \item \textbf{Field solving}: variational relaxation for 
          high-dimensional PDEs.
\end{itemize}
These problems either lack or are not well-served by exact 
polynomial algorithms. This is where \ZSP's iterative nature becomes 
an \emph{advantage}---the physics itself performs the iteration in 
parallel, at speeds no digital computer can match.

\subsection{Why the Negative Result Matters}

Negative results are underexplored in machine learning and 
computing research. They have unique value:
\begin{enumerate}[leftmargin=*]
    \item \textbf{Boundary mapping}: they delineate what a method 
          cannot do, saving future researchers from repeating 
          the same dead ends.
    \item \textbf{Theory validation}: they confirm theoretical 
          predictions (here, LP duality) in practice.
    \item \textbf{Resource reallocation}: they redirect effort 
          toward problems where the method has an advantage.
\end{enumerate}

\subsection{Future Directions}

Based on our analysis, we recommend:
\begin{enumerate}[leftmargin=*]
    \item \textbf{Material simulation}: implement the Hadamard 
          grid for entanglement fingerprinting.
    \item \textbf{Physics-layer cryptography}: prototype the 
          PUF / limit-cycle key agreement.
    \item \textbf{FPGA mapping}: convert the digital twin to 
          Verilog-A / SPICE netlists.
    \item \textbf{Silicon photonics tape-out}: partner with 
          a foundry for 40nm CMOS + SOI MPW.
\end{enumerate}
\section{Conclusion}
\label{sec:conclusion}

We systematically evaluated Zero-Space Projection (\ZSP) on three 
perception tasks: single-frame \MOT, multi-frame \MOT, and 
cooperative perception. Our conclusions are:

\begin{enumerate}[leftmargin=*]
    \item On \textbf{single-frame \MOT}, \ZSP achieves $99.96\%$ of 
          Hungarian's accuracy but is $20{,}000\times$ slower.
    \item On \textbf{multi-frame \MOT}, \ZSP is statistically 
          significantly worse than Hungarian ($-6.7\%$, $p=0.0001$, 
          win rate $0\%$).
    \item On \textbf{cooperative perception}, all methods achieve 
          $100\%$ accuracy, but distributed Hungarian is fastest and 
          requires zero communication.
\end{enumerate}

\textbf{Root cause.} Linear assignment problems possess a dual 
structure that enables the Hungarian algorithm to achieve exact 
optimality in polynomial time. \ZSP is an iterative approximation 
and cannot, by construction, outperform the exact algorithm.

\textbf{Implication.} The value of \ZSP---and of physics-native 
computing more broadly---is not in replacing problems already 
solved by exact algorithms. It is in solving problems that lack 
such algorithms: material simulation, physics-layer cryptography, 
high-precision sensing, and variational field solving.

\textbf{Final thought.} Knowing when \emph{not} to use a tool is 
as important as knowing when to use it. We hope this negative 
result serves as a useful boundary marker for the physics-native 
computing community.

\bibliographystyle{plain}
\bibliography{refs}

\newpage
\appendix
\section{Experimental Configuration}
\label{app:config}

\subsection{Hardware and Software}

All experiments were run on:
\begin{itemize}[leftmargin=*]
    \item CPU: Intel Core i9-13900K (24 cores, 5.8 GHz)
    \item RAM: 64 GB DDR5
    \item OS: Windows 11
    \item Python: 3.12.10
    \item NumPy: 2.x
    \item SciPy: 1.14.x
    \item pytest: 9.0.1
\end{itemize}

\subsection{Test Configuration Table}

\begin{table}[H]
    \centering
    \caption{Complete experimental matrix.}
    \label{tab:config}
    \begin{tabular}{l l l l}
        \toprule
        Experiment & Config & Method & Metrics \\
        \midrule
        Single-frame & $T,D \in \{10,20,30,50\}$ & Hung./Greedy/ZSP & Cost, Time \\
        Multi-frame & $K \in \{3,4,5\}$ & Hung./ZSP & Acc., IDSW \\
        Cooperative & $M \in \{2,3,4\}$ & Cent./Dist./ZSP & Acc., Time, Comm \\
        Statistical & 30 seeds & Paired $t$-test & $p$-value \\
        \bottomrule
    \end{tabular}
\end{table}
\section{Convergence Analysis}
\label{app:convergence}

\begin{proposition}[Monotone Descent]
Let $J: \R^K \to \R$ be $L$-smooth and convex. The projected gradient 
update
\begin{equation}
    \avec_{k+1} = \avec_k - \eta \nabla J(\avec_k)
\end{equation}
satisfies
\begin{equation}
    J(\avec_{k+1}) \le J(\avec_k) - \eta \left(1 - \frac{\eta L}{2}\right) 
    \|\nabla J(\avec_k)\|^2.
\end{equation}
\end{proposition}

\begin{proof}
By $L$-smoothness:
\begin{equation}
    J(\avec_{k+1}) \le J(\avec_k) + \nabla J(\avec_k)^\top 
    (\avec_{k+1} - \avec_k) + \frac{L}{2} \|\avec_{k+1} - \avec_k\|^2.
\end{equation}
Substituting $\avec_{k+1} - \avec_k = -\eta \nabla J(\avec_k)$:
\begin{align}
    J(\avec_{k+1}) &\le J(\avec_k) - \eta \|\nabla J\|^2 
    + \frac{\eta^2 L}{2} \|\nabla J\|^2 \\
    &= J(\avec_k) - \eta \left(1 - \frac{\eta L}{2}\right) \|\nabla J\|^2.
\end{align}
For $\eta < 2/L$, the right-hand side decreases monotonically.
\end{proof}

\begin{corollary}[Convergence Rate]
For $\eta < 2/L$, the iterates converge to the global optimum 
$\avec^*$ with rate:
\begin{equation}
    J(\avec_k) - J(\avec^*) \le \frac{\|\avec_0 - \avec^*\|^2}{2\eta k}.
\end{equation}
\end{corollary}

\subsection{Non-Negative Projection}

After each step, we clamp $\xvec \ge 0$. This is equivalent to 
projecting onto the non-negative orthant, which is a convex set. 
By standard projected gradient theory~\cite{bertsekas1999nonlinear}, 
convergence is preserved.
\section{Code and Data}
\label{app:code}

\subsection{Repository Structure}

\begin{verbatim}
abn_qss_mot/
├── core/
│   ├── constraint.py       # Constraint matrix construction
│   ├── nullspace.py        # SVD-based null-space extraction
│   ├── solver.py           # Projected gradient solver
│   ├── multi_frame.py      # Multi-frame extension
│   └── cooperative.py      # Cooperative extension
├── experiments/
│   ├── benchmark.py        # Single-frame benchmark
│   ├── benchmark_multiframe.py
│   ├── cooperative_benchmark.py
│   ├── statistical_significance.py
│   └── plotting/
│       ├── fig1_digital_twin.py
│       └── ... (8 figures)
├── tests/
│   ├── test_constraint.py
│   ├── test_nullspace.py
│   └── test_solver.py
└── results/
    ├── single_frame.txt
    ├── multi_frame.txt
    └── statistical.txt
\end{verbatim}

\subsection{Reproduction Commands}

\begin{verbatim}
# Install dependencies
pip install -e .

# Run all tests
pytest tests/ -v --cov=core

# Single-frame benchmark
python -m experiments.benchmark

# Multi-frame benchmark
python -m experiments.benchmark_multiframe

# Cooperative benchmark
python -m experiments.cooperative_benchmark

# Statistical significance
python -m experiments.statistical_significance

# Regenerate all figures
python -m experiments.plotting.plot_all
\end{verbatim}

\subsection{License}

Code released under MIT license. Data released under CC-BY-4.0.

\end{document}